% LSPV-3 REPORT::DEPDEF1::2026-10-01T11:36:24-07:00
% !TEX program = xelatex
\documentclass[11pt]{article}
\usepackage[margin=0.88in]{geometry}
\usepackage{amsmath,amssymb}
\begin{document}
\section*{LSPV-3: dependency-preserving reconstruction audit}
\textbf{STATUS: OPEN.} P versus NP remains open. This report sharpens the dependency obstruction in the reconstruction collapse. It proves no unrestricted SAT lower bound.
\paragraph{CORPUS REVIEWED.} I surveyed the current repository protocol, root living main, current and historical MAIN snapshots surfaced by repository search, active task specifications, materially relevant LSPV-1--5 reports, the random-challenge HEC/CRS--PPC branch and theorem index, governance material surfaced by search, and the experiment/result artifacts cited by the active proof spine. In particular I checked HADCOEF1, FEAT1/CODE1, RECDEP1/QREC1/PQCOL1, ACQ1, the current 10:46 MAIN snapshot, HEC-SAT proved core, CRS--PPC theorem index, and TASK-004. Predecessor LSPV-2 FEAT1 is substantive, has the required header and sentinel, and reads back from GitHub; preflight passes.
\paragraph{Two obstacles.} First, the final universal block in the reconstruction predicate ranges over assignments to SAT instances whose descriptions are chosen after the middle existential variables. Second, any proposed uniform certificate must be acquired, represented, verified, compiled, and written in the same realization; existence or definability alone is not a quantitative collapse.
\paragraph{Primary-literature preflight.} Viola proves $\mathrm{BPTime}(t)\subseteq\Sigma_3\mathrm{Time}(t\,\mathrm{polylog}\,t)$ and explicitly states that this simulation relativizes. van Melkebeek--Mocelin Sdroievski provide the resilient reconstruction architecture. Recent DQBF work makes explicit that existential variables may have prescribed dependency sets and that unrestricted DQBF satisfiability is NEXP-complete. Definition-based QBF dependency schemes can detect additional independencies and yield exponential proof savings on some families, but they are proof-system/formula-structure results, not unrestricted SAT-time quantifier elimination.
\paragraph{PROVED/EXACT: bundled reconstruction semantics.} Fix one acquired, size-normalized SAME\_REALIZATION $Q$, coordinate $i$, and Arthur coins $r$. For the constant family $J$ of generated SAT instances, with verifier $\Phi_j$, exact answers satisfy
\[\operatorname{ExactAns}(a;Q,i,r)=\exists Y^+\,\forall Y^-\bigwedge_{j\in J}[(a_j\Rightarrow\Phi_j(Y_j^+))\wedge(\neg a_j\Rightarrow\neg\Phi_j(Y_j^-))].\]
Bundling same-polarity witnesses is exact but does not remove an invocation of the assumed SAT decider.
\paragraph{PROVED/EXACT: dependency graph obstruction.} After late substitution into the post-Viola simulation, the project has the schematic prefix
\[\boxed{\exists\sigma\;\forall u\;\exists z\;\forall y\;R(x,Q,\sigma,u,z,y)},\qquad z=(\ell,a,Y^+),\ y=Y^-.\]
Let $\Phi_{u,z}$ denote the downstream SAT instance whose NO semantics are checked by $y$. The universal assignment domain is therefore indexed by $(u,z)$ through the description of $\Phi_{u,z}$. A sound commutation of $\forall y$ before $\exists z$ is justified only if one supplies a representation $Z$ independent of $z$ together with a deterministic predicate $V$ such that, for every legal $z$, $V(x,Q,\sigma,u,Z,z)$ certifies exactly the required NO facts for $\Phi_{u,z}$. This is a dependency-uniformization condition, not ordinary witness bundling.
\paragraph{PROVED/EXACT: no free flattening from definability.} Even if such a function $Z=F(x,Q,\sigma,u)$ is mathematically unique or definable, that alone does not improve the exponent ledger. The run must charge construction/acquisition of $Z$, its length, indexing into it, verification, compilation, and output writing. Thus a definition-based dependency reduction is quantitatively useful here only with an explicit model-matched algorithm whose total exponent is smaller than the eliminated recursive SAT layer.
\paragraph{KNOWN THEOREM / transfer boundary.} DQBF is the exact formal language for explicitly restricted existential dependencies, and current fine-grained DQBF work studies how dependency sets alter complexity. Definition-based dependency schemes are a genuine proof-complexity mechanism for proving some variables independent of others. Their transfer to this reconstruction is \textbf{ANALOGY/CANDIDATE}: we still need a theorem proving the required independence for the concrete vMS/Viola predicate and an acquisition algorithm. Treating a dependency scheme itself as an unrestricted RAM lower bound is a \textbf{FAILED TRANSFER}.
\paragraph{Exponent ledger.} If the fixed exact SAT decider runs in $S^{k+o(1)}$ and the four actual model-matched compilation stages have exponents $b_0,b_1,b_2,b_3$, the conservative recurrence remains
\[\boxed{\alpha_{\rm det}\le \alpha_0(b_0b_1b_2b_3)k^4+o(1).}\]
This is \textbf{CONDITIONAL} upper-bound accounting, not a lower bound and not a theorem that $k^4$ is necessary. A genuine $k^3$ improvement requires eliminating one charged invocation layer or replacing it by a strictly cheaper structured compiler/uniformizer.
\paragraph{Separate cost ledger.} The same realization must charge
\[\tau_Q+\tau_{\rm norm}+\tau_{\rm census}+\tau_{\rm amp}+\tau_{\rm idx}+\tau_{\rm compile}+\tau_{\rm rec}+\tau_{\rm write}+\tau_{\rm transport}.\]
TransportCert is only an accounting interface. ACQ1 independently shows that common-$Q$ acquisition can itself require charged SAT calls when implemented by legal-state separation.
\paragraph{PQ structural toolkit.} \textbf{EXACT OBJECT/THEOREM TRANSFER:} Schur squares, shortening/puncturing, quadratic-relation kernels, fixed-encoder slice matrices, and collision Cauchy--Schwarz on their declared maps. \textbf{ANALOGY/CANDIDATE:} Hilbert/syzygy invariants may expose algebraic dependencies; special-soundness style collision extraction could help only if two conflicting valid vMS openings yield an explicit locally checkable forbidden witness; an LWE-style black-box architecture could help only via a new construction-specific external extraction theorem. \textbf{FAILED TRANSFER:} cryptographic collapse-binding, LWE hardness, random oracles, VLSI locality, control barriers, distributed cuts, and proof-system dependency savings do not imply unrestricted SAT runtime.
\paragraph{SAME\_REALIZATION, TC4, relativization, algebrization.} The proposed $Z$ must belong to the same $Q$, Arthur randomness, encoder, legal-state family, and downstream query family. TC4 remains OPEN because an arbitrary exact SAT algorithm may bypass representation-specific objects unless an external reduction forces or extracts them. Viola's simulation and the recursive SAT-elimination accounting relativize; dependency bookkeeping and the present algebraic code objects do not supply a nonrelativizing hinge. The algebraic pieces are also visible to algebraic oracles, so no nonalgebrizing step has been identified.
\paragraph{DUPLICATION / CONTRADICTION AUDIT.} RECDEP1 already identified uniformization as the only plausible local flattening. The new contribution is to state its exact dependency condition and to reject a possible overreading suggested by modern dependency-scheme literature: proving definability or proof-system independence is insufficient without a charged uniformizer. No contradiction with FEAT1, HADCOEF1, ACQ1, HEC, or CRS--PPC was found. The causal random-challenge branch remains orthogonal and restricted-model until its static compiler closes.
\paragraph{WHAT CHANGED RELATIVE TO ROOT MAIN.} Root main still cites the older 2026-09-30 reconstruction report and states only the raw four-block obstruction. It is stale with respect to QREC1/RECDEP1. This run updates the specialist state to include the exact recursive recurrence, one-$k$ YES/NO accounting, and the dependency-uniformization criterion while preserving OPEN status.
\paragraph{ADOPTED / REJECTED / RETIRED PRIOR COMPONENTS.} Adopted: the shared-$Q$ exact-answer predicate, eager-oracle factorization, recursive $b_qk$ recurrence, ACQ1 acquisition charge, and strict fiber/collision formulation. Rejected as insufficient: generic canonical/unique binding and uncharged definability. Retired: counting YES and NO as separate $k$ factors; qualitative $P=NP\Rightarrow PH=P$ as a quantitative compiler; direct cryptographic or physical transfer.
\paragraph{STALE OR OVERSTATED CLAIMS FOUND.} The LSPV-3 owned region of root main is stale because it points to the older 42301-R1 report and omits the later exact recurrence and uniformization criterion. No verified theorem was found that justifies $k^3$.
\paragraph{Strict status.} \textbf{PROVED/EXACT:} exact-answer predicate; dependency-uniformization condition as a necessary condition for the proposed commutation; definability alone does not erase charged acquisition/verification; recursive elimination accounting under declared compilers. \textbf{KNOWN THEOREM:} Viola simulation and relativization; DQBF explicit dependency semantics. \textbf{CONDITIONAL:} four-stage $k^4$ upper-bound ledger. \textbf{CANDIDATE:} construction-specific uniformizer, collision extractor, Hilbert/syzygy diagnostics, LWE-style external extraction. \textbf{FALSIFIED/RETIRED:} free canonical binding, free definability, YES/NO double exponent, qualitative PH collapse as quantitative compiler. \textbf{OPEN:} existence and cost of a vMS/Viola uniformizer; exact compiler exponents; full $Q$ acquisition/normalization; TC4; nonrelativizing/nonalgebrizing hinge; final diagonal.
\paragraph{Single next target.} Write the actual post-Viola bounded predicate with a dependency table mapping each downstream SAT-instance description bit to $(Q,r,u,z)$. Then test whether the NO-side dependency on $z$ is essential by finding either (i) a construction-specific $Z=F(Q,r,u)$ plus a cheaper verifier, or (ii) two legal middle choices $z_0,z_1$ requiring incompatible NO certificates under the same upstream state. Only (i), with fully charged construction, can justify recomputing a $k^3$ ledger.
\begin{thebibliography}{5}
\bibitem{Viola05} E. Viola, \emph{On Probabilistic Time versus Alternating Time}, ECCC TR05-137, 2005.
\bibitem{vMS25} D. van Melkebeek and N. Mocelin Sdroievski, \emph{Instance-Wise Hardness and Refutation versus Derandomization for Arthur--Merlin Protocols}, Computational Complexity 34:15, 2025.
\bibitem{Cheng25} C. Cheng, L.-H. Fung, J.-H. R. Jiang, F. Slivovsky, T. Tan, \emph{Fine-Grained Complexity Analysis of Dependency Quantified Boolean Formulas}, SAT 2025, LIPIcs 341:10.
\bibitem{DefDep26} \emph{Definition-Based Dependency Schemes}, SAT 2026, LIPIcs, 2026.
\bibitem{AW09} S. Aaronson and A. Wigderson, \emph{Algebrization: A New Barrier in Complexity Theory}, ACM TOCT 1(1), 2009.
\end{thebibliography}
\end{document}
% END_OF_VERIFIED_REPORT::LSPV-3::DEPDEF1